% Figure: Force-Voltage Analogy
% Chapter 5: Electromechanical Analogies
\begin{figure}[htbp]
\centering
\begin{tikzpicture}
    % ========== MECHANICAL SYSTEM (Left) ==========
    \node[font=\bfseries, UofSCGarnet] at (0, 4) {Mechanical System};

    % Wall
    \fill[pattern=north east lines] (-0.3, 0) rectangle (0, 3);
    \draw[thick] (0, 0) -- (0, 3);

    % Spring (top)
    \draw[thick] (0, 2.5) -- (0.3, 2.5);
    \draw[thick, decorate, decoration={coil, segment length=4pt, amplitude=4pt}]
        (0.3, 2.5) -- (2.0, 2.5);
    \draw[thick] (2.0, 2.5) -- (2.3, 2.5);
    \node[above, font=\small] at (1.15, 2.7) {$k$};

    % Damper (middle)
    \draw[thick] (0, 1.5) -- (0.4, 1.5);
    \draw[thick] (0.4, 1.2) rectangle (1.2, 1.8);
    \draw[thick] (0.7, 1.3) -- (0.7, 1.7);
    \draw[thick] (1.2, 1.5) -- (2.3, 1.5);
    \node[below, font=\small] at (0.8, 1.1) {$b$};

    % Mass
    \draw[thick, fill=UofSCAtlantic!20] (2.3, 0.8) rectangle (3.5, 2.8);
    \node[font=\bfseries] at (2.9, 1.8) {$m$};

    % Force arrow
    \draw[->, thick, UofSCRose, line width=2pt] (3.7, 1.8) -- (4.5, 1.8)
        node[right, UofSCRose, font=\small] {$F(t)$};

    % Velocity arrow
    \draw[->, thick, UofSCCongaree] (2.9, 0.5) -- (3.5, 0.5)
        node[right, font=\small, UofSCCongaree] {$v(t)$};

    % Ground
    \draw[thick] (-0.3, 0) -- (4.8, 0);
    \foreach \x in {0, 0.4, ..., 4.4} {
        \draw[thick] (\x, 0) -- (\x-0.15, -0.25);
    }

    % Equation
    \node[font=\small, align=center] at (2.2, -0.8) {
        $F = m\dfrac{dv}{dt} + bv + k\displaystyle\int v\, dt$
    };

    % ========== ANALOGY ARROWS (Center) ==========
    \draw[<->, very thick, UofSCGarnet] (5.2, 2.5) -- (6.8, 2.5);
    \node[font=\bfseries\small, UofSCGarnet] at (6, 3) {Analogy};

    % Analogy table
    \node[draw=UofSC70Black, fill=UofSCSandstorm, font=\footnotesize,
          align=center, sharp corners, inner sep=4pt] at (6, 1.2) {
        $F \leftrightarrow V$\\[2pt]
        $v \leftrightarrow i$\\[2pt]
        $m \leftrightarrow L$\\[2pt]
        $b \leftrightarrow R$\\[2pt]
        $k \leftrightarrow 1/C$
    };

    % ========== ELECTRICAL SYSTEM (Right) ==========
    \node[font=\bfseries, UofSCGarnet] at (10.5, 4) {Electrical System (RLC)};

    \begin{scope}[shift={(7.5, 0.5)}]
        \begin{circuitikz}[scale=0.9]
            % Voltage source
            \draw (0, 0) to[V, l=$V(t)$, invert] (0, 3);

            % Inductor
            \draw (0, 3) to[L, l=$L$, i=$i$] (2, 3);

            % Resistor
            \draw (2, 3) to[R, l=$R$] (4, 3);

            % Capacitor
            \draw (4, 3) to[C, l=$C$] (4, 0);

            % Ground return
            \draw (4, 0) -- (0, 0);

            % Current arrow label
            \node[UofSCCongaree, font=\small] at (3, 3.6) {$i(t)$};
        \end{circuitikz}
    \end{scope}

    % Equation
    \node[font=\small, align=center] at (10.5, -0.8) {
        $V = L\dfrac{di}{dt} + Ri + \dfrac{1}{C}\displaystyle\int i\, dt$
    };

    % ========== Element Labels ==========
    \node[UofSCAtlantic, font=\footnotesize] at (2.9, 3.3) {Inertia};
    \node[UofSCAtlantic, font=\footnotesize] at (8.8, 3.8) {Inertia};

    \node[UofSCHorseshoe, font=\footnotesize] at (0.8, 0.6) {Dissipation};
    \node[UofSCHorseshoe, font=\footnotesize] at (10.5, 3.8) {Dissipation};

    \node[UofSCCongaree, font=\footnotesize] at (1.15, 3.2) {Storage};
    \node[UofSCCongaree, font=\footnotesize] at (12.2, 1.5) {Storage};

\end{tikzpicture}
\caption{Force-voltage analogy between mechanical and electrical systems. The mass-spring-damper system and series RLC circuit are governed by mathematically identical differential equations, enabling cross-domain analysis techniques.}
\label{fig:force_voltage_analogy}
\end{figure}
